\ssscp{Medial consonant}{14-04-03-03}
\citet[242]{Blust1982} discusses medial *ns in **mans(aə)r
and finds that it is not distinguished from *nd in eastern Indonesia,
but concludes that the reflexes are consistent with *ns.
Whatever the evidence from EMP,
within our target region,
with the exception of certain \tsc{Ambon-Seram} languages that
clearly reflect *ŋinsi rather than *ŋisi ‘teeth’,
data unambiguously reflecting *ns are so scarce as to forestall
making any claims of regular sound correspondences with any confidence.
But because of the merger of *nt/*nd/*ns > **nd in \tsc{Ambon-Seram}
(see below), *nt and *nd are equally valid,
with more data available on which to verify regular sound correspondences.

Medial NC clusters are found only in the Sulawesi forms in the
western part of the range \it{mante},
and in Yamdena \it{manr-e} [manre],
[mandre], [mande], Kei \it{mendar}, and Banda \it{mender,
mbender} in the east.
No known language in our target region preserves surface \it{ns}.
Other languages have reduced to a single medial consonant (see \srf{13-02-01}).

Beginning with the Sulawesi forms,
medial \it{nt} in \it{mante} would not be regular from *ns.
\citet[38--40]{MeadTruong2021} report *ns = \it{ns} in Lasalimu,
but *ns > \it{nc} [ɲʧ] in all other languages of Buton.
They further report *nt = \it{nt} in Kaisabu,
with a split of *nt > \it{nt,
nʧ} in other languages of Buton.
An example of *ns > \it{ns,
nʧ} is *ŋusuq ‘nasal area’ > **ŋunsu ‘labial circle’ > Lasalimu
\it{ŋunsu}, Wasambua, Wakole, Wolowa, Kumbewaha \it{ŋunʧu}.
An example of *nt = \it{nt} is *muntay ‘citrus’ > all Buton \it{munte}.
See \citet[38--40]{MeadTruong2021} for more examples and discussion.
Thus, \tsc{Muna-Buton} \it{mante} ‘cuscus’ points to earlier
medial *nt, and rules out earlier *ns.

\citet[41]{Stresemann1927} and \citet[36]{Collins1981}, \citeyear[55]{Collins1983} note
that in Central Maluku languages *nt,
*nd, and *ns merge
and have the same reflex.
For \citet{Collins1986}, “Central Maluku” includes the \tsc{East
Seram} and \tsc{Eastern Islands} nodes of our \tsc{Seram-Tanimbar-Bomberai} subgroup.
\citet[43, endnote 29]{Collins1981} notes, “Evidence from Alune confirms
the reconstruction of *ns in PCM [**manser]”.
Seven varieties of Alune show \it{marel(e)}.
One variety (Kairatu) has \it{maʧel}.
But limiting the medial consonant to *ns is too restrictive,
since *nt is equally likely
and is more robustly attested.
For example: PMP *punti ‘banana’ > Alune \it{uri} (Riring
and Sapolewa varieties), \it{huri} (Laturake variety), and \it{uʧi} (Kairatu).
This shows that PMP *nt > \it{r} in most Alune varieties (> \it{ʧ}
in Kairatu), so the Alune forms can derive from **mantəR ‘cuscus’ through regular sound correspondences.
The same is true for North Wemale (\it{marel(e)} ‘cuscus’, cf.
\it{huri} ‘banana’) and South Wemale (\it{madel(e)} ‘cuscus’, cf.
\it{hudi} ‘banana’), and many other languages.
Likewise, the \tsc{Piru Bay} languages of \tsc{Ambon-Seram} with
medial \it{k}, such as Paulohi \it{makel-e} ‘cuscus’,
can also derive from medial *nt (cf.
*punti > Paulohi \it{fuki}).

The word \it{mate} ‘cuscus’ in the \tsc{Sula-Buru} language Ambelau
could be a loan from a southeast Sulawesi language,
or a retention.\footnote{
Given that Buru \it{tonal} ‘cuscus (generic)’
is a loan from a southeast Sulawesi source (cf.
Kambowa \it{tonali} ‘dwarf cuscus,
\it{Strigocuscus celebensis}’), the Ambelau form could also be
a loan from \it{mante},
given the large presence of thousands of southeast Sulawesi migrants
along coastal south and west Buru, particularly from Cia-Cia.}
Ambelau words with
a medial surface \it{t} are rare in the available data due to
regular *t > \it{r} (\srf{09-02}),
but we note that in the one available form *nt > \it{t}; **ntibu
> \it{-tiβo} ‘fly’ (see example \qf{07-03}, \crf{ch:CM}).
More broadly, \citeauthor{Stresemann1927}’s (\citeyear[41]{Stresemann1927})
and \citeauthor{Collins1981}’ (\citeyear[36]{Collins1981}; \citeyear[55]{Collins1983}) assertion of a merger in central
Maluku of *nt/*nd/*ns > **nd is problematic for \tsc{Sula-Buru}.
For Buru, the scarce data available suggest that *nd > \it{d},
*nt > \it{t}, *ns > \it{s},
showing a consistent loss of the nasal
and preservation of the following consonant.
(Normally for Buru, *d > \it{r},
*t = \it{t}, *s = \it{s}.) So the claimed merger for “Central
Maluku” of *nt/*nd/*ns > **nd is incorrect for \tsc{Sula-Buru}.
If Ambelau patterns like Buru,
and \it{mate} ‘cuscus’ is a retention rather than a loan,
then the \tsc{Sula-Buru} data exclude medial *ns or *nd,
but are consistent with medial *nt.

\tsc{Seram-Tanimbar-Bomberai} data are divided among several etyma meaning ‘cuscus’.
Most languages in the \tsc{Eastern Islands} node are accounted for
by **kVndoR(a) ‘cuscus’ (\srf{14-03-04}),
and not **mantəR ‘cuscus’.
The exception is Bati \it{manuk faniki} `cuscus' from *manuk `bird' and *paniki `bat'.
Similarly \tsc{East Seram} Masiwang also derives
its word for ‘cuscus’ from PMP *manuk ‘chicken, bird’.
The Bati form \it{manuk faniki} ‘cuscus’,
further links the word for ‘cuscus’ with the PMP etymon for
*paniki ‘fruit bat’, another edible animal found in the treetops
(\citealp[269f]{Schapper2011}, see \srf{14-04-04}).

Data in the \tsc{Tanimbar-Bomberai} node reflect **mandəR or
**məndəR, or a form with medial *nt.
Medial *ns is difficult to support.
The medial consonants of Yamdena \it{manr-e} Kei Besar \it{mendar},
Kei Kecil \it{medar}, Fordata \it{medar} would be expected from
a nasal + \it{d}.
Consider PMP *də\tbr{md}əm ‘dark’ > Yamdena \it{nro\tbr{nr}am}
‘night’, Kei Kecil \it{de\tbr{d}an} ‘entire night’,
Fordata \it{de\tbr{d}ən} ‘pitch black’; **du\tbr{md}um ‘thunder
rumbling’ > Yamdena \it{nru\tbr{nr}um} ‘rumbling thunder’,
Kei Kecil \it{do\tbr{d}} ‘thunder’; PMP *ma-di\tbr{ŋd}iŋ ‘cold’
> Kei Kecil \it{na-ri\tbr{d}in} ‘cold’,
Fordata \it{ri\tbr{d}in} ‘cold, cool’,
(but Yamdena \it{-ndi\tbr{r}in} ‘cold’).
There are other correlations between Yamdena \mbox{\it{-nr-}}
and Fordata \it{-d-}, such as Yamdena \it{-si\tbr{nr}ak} ‘split’
and Fordata \it{n-si\tbr{d}aŋ} ‘split roughly’.
The possibility of medial *nt can also be countenanced with evidence
of post-nasal voicing, as in *titis ‘drip,
ooze’ > Kei Kecil \it{n-\tbr{d}it} ‘drip’, Fordata \it{n-\tbr{d}iti} ‘drip’.

Banda \it{mender} ‘cuscus’ \citep[196]{Ellen1993}
and \it{mbender} ‘cuscus’ \citep[45]{Kakerissa1986} could reflect a medial *nt (cf.
PCM **ntibu ‘fly (v.)’ > Banda \it{tifu {\tl} ndifu} [inflected
variants]); medial \it{t} (cf.
Banda \it{ni-ka\tbr{nd}a} ‘\tsc{3pl.gen}-trousers’, Asilulu \it{ri-ka\tbr{t}a}
‘\tsc{3pl.gen}-trousers’, Buru \it{nun(u) ka\tbr{t}a-ro} ‘\tsc{3pl.gen} trousers-\tsc{pl}’) or *nd (cf.
PMP *ma-di\tbr{ŋd}iŋ ‘cold’ > Banda \it{ndiri\tbr{nd}in} ‘feel cold’).

Proto-Aru **medaR ‘cuscus’ \citep[264]{Schapper2011} which would
be regular from **məntəR,
accounts for the cognate data from Aru languages.
Historical *nt > **d in Aru (Rick Nivens, pers. comm. March 2021).

The medial consonants of \tsc{Timor-Babar} are mostly accounted
for regularly by *nt and/or **nd.
Ili{\Q}uun (Wetar Island) \it{maʧa} has regular **nd > \it{ʤ {\tl} ʧ}.
Kisar \it{masla} appears to present some challenges.
Given that PMP *s > \it{h} in Kisar,
this looks like \it{masla} might be evidence for a medial *ns.
But consider the data in \trf{14-10} which shows Kisar \it{s}
corresponding to Roma and Leti \it{d}, Luang \it{t}.
The reflexes of *paŋdan ‘pandanus’ indicate that all these are from earlier *nd.
(Irregular initial *a > \it{e} in *paŋdan is also attested in
the \tsc{Rote-Meto} node of \tsc{Timor-Babar} [\cite[173]{Edwards2021}].)

\begin{table}[h]
	\centering\tcp{\tsc{Kisar-Luang} **nd}{14-10}
	\st{6pt}\begin{tabular}{l|l@{ }l@{ }ll}\lsptoprule
gloss						&	\mc{2}{l}{{\hr}pandanus}							&	rattan														&	kill\\
source					&	\mc{2}{l}{\hr*pa\tbr{ŋd}an}						&	?Malay \it{rotan}									&	?\\	\midrule
P.-Kisar-Luang	&	\mc{2}{l}{**e\tbr{nd}an}							&	**lo\tbr{nd}an										&	**we\tbr{nd}an\\
Kisar						&	{\hr\hr}e\tbr{s}ne	&	‘pineapple’			&	{\hr\hr}lo\tbr{s}ne, lo\tbr{s}no	&	{\hr\hr}n-e\tbr{s}ne\\
Roma						&	{\hr\hr}e\tbr{d}na	&	‘pandanus’			&	{\hr\hr}lo\tbr{t}na								&	{\hr\hr}n-e\tbr{d}ina\\
Luang						&	{\hr\hr}e\tbr{t}nə 	&	‘pandanus leaf’	&	{\hr\hr}lo\tbr{t}nə								&	{\hr\hr}n-we\tbr{n}nə\\
Leti						&	{\hr\hr}ɛ\tbr{d}an	&	‘pineapple’			&	{\hr\hr}lɔ\tbr{d}an								&	{\hr\hr}n-βe\tbr{d}na, n-βenna\\
		\lspbottomrule
	\end{tabular}
\end{table}

In the \tsc{Babar} node, Southeast Babar and East Marsela \it{mɔt},
Tela \mbox{\it{mut-ɛ}}, and North Babar \it{mad-a} are accounted for by *nt,
as shown by *punti > Southeast Babar, East Marsela, Tela \it{uty} `banana',
North Babar *punti + **biak {\ra} \it{udwiya}.
Emplawas \it{mutt-o} ‘cuscus, civet’ with double \it{tt} may present a challenge.
The only (other clear) instance of *nt in Emplawas is *punti > \it{uty-o}.\footnote{
		One possible instance of *nd in Emplawas is
		*tindəs > \it{-tret} `crush, push down'.}
While Emplawas \it{mutt-o} may be evidence for a consonant other than *nt/*nd,
a local development cannot be ruled out.
It is possible that the medial consonant is from intermediate
**ntR/**dR after loss of the final vowel,
or it could be conditioned by the previous vowel.\footnote{
		\tsc{South Babar} languages have *a > **u /{\#}(N){\gap} for which \citet{Zobel2018}
		posits intermediate **ə. Schwa often has certain unique properties in our target region,
		including lengthening of a following consonant (see \srf{03-02-01-02} and \srf{03-03-01-02}).}

Makuva \it{mate} could attest *s > \it{t},
but *nt > \it{t} also occurs,
as in PMP *punti > \it{utu} ‘banana’.
The forms in \tsc{Central Timor} languages are regularly accounted
for by medial *nt or *nd.
Mambae \it{mat, maat} ‘cuscus’ may point to earlier *nt rather than *nd.

To summarise, \tsc{Muna-Buton} languages exclude *ns and *nd.
If Ambelau \it{mate} ‘cuscus’ is a retention,
the \tsc{Sula-Buru} data exclude both a medial *ns and *nd.
Either a medial *nt or *nd are supported by several Wallacean subgroups.
Thus, we conclude that the medial consonant was most likely *nt,
with subsequent parallel *nt > **nd in many subgroups through
the common process of post-nasal voicing.

